% Introduction

\subsection{Motivation}

Cellular decision-making requires integrated models capturing metabolic and regulatory information flow. Biological systems exhibit hierarchical organization where metabolic state (e.g., energy status) controls regulatory networks (e.g., gene expression), which in turn modulate metabolic flux. This bidirectional coupling demands unified formalisms treating metabolism and regulation as inseparable components of cellular control, not separate modeling domains. For instance, ATP in \textit{Bacillus subtilis} simultaneously fuels metabolism (mass transfer) and gates developmental commitment (regulatory signal)---when ATP falls below 2.21 mM, cells irreversibly commit to sporulation, a threshold current formalisms cannot compute.

\noindent Biological Petri nets~\cite{reddy1993,heiner2008} extended classical formalism with continuous places for metabolites, stochastic semantics for molecular noise, and test arcs for catalysis. However, a critical limitation remains: metabolism and regulation are separated. Normal arcs model mass transfer (substrates consumed, products generated), while test arcs model catalytic reads (enzymes unchanged). This dichotomy cannot capture regulatory molecules \textit{consumed} during decision-making---signal consumption that collapses decision spaces and ensures commitment finality. Growth factor binding depletes extracellular signal, ATP expenditure in \textit{B. subtilis} sporulation gates developmental commitment, and similar consumption occurs in lambda phage CI repressor binding. Recent advances address resource constraints~\cite{genovese2021}, but no framework provides formal hierarchical information channels with consumption semantics necessary for modeling integrated metabolic-regulatory control.

\subsection{Biological Motivation for Signal Hierarchy}

Cellular regulation exhibits multi-scale hierarchical organization: \textbf{Stress response}---Environmental signals (UV damage, heat shock) trigger decision points that reconfigure metabolic pathways through hierarchical preemption where regulatory layers control metabolic execution. \textbf{Energy status}---ATP/ADP ratio serves as a master regulatory signal, hierarchically controlling anabolic versus catabolic flux through multiple pathways simultaneously. \textbf{Developmental decisions}---Nutrient depletion in \textit{Bacillus subtilis} triggers ATP-gated commitment to sporulation versus competence pathways, exemplifying metabolic state controlling gene regulatory networks. \textbf{Quorum sensing}---Autoinducer concentration represents population-level information that preempts individual cell decision-making, enabling coordinated multicellular behavior.

\noindent 
Existing Bio-PNs model catalysis via test arcs (non-consuming reads). However, regulatory signals are often \textit{consumed} during decision-making, changing system state irreversibly. For example, ATP consumption in \textit{B. subtilis} sporulation creates distinct attractor basins for developmental commitment, and growth factor binding to receptors depletes extracellular signal. Lambda phage provides another example where CI repressor binding to operators consumes regulatory signal concentration. We formalize signal places ($\Psi$) with signal flow arcs that consume tokens, enabling hierarchical preemption mechanisms absent in current Bio-PN theory.

\subsection{Contributions}

This work makes two primary contributions:

\noindent\textbf{(1) Signal Hierarchy Theory}---We introduced signal places ($\Psi$) as formal information channels distinct from metabolic mass transfer. Signal flow arcs consume tokens (unlike test arcs which read without consumption), enabling hierarchical preemption where higher regulatory layers control lower metabolic layers. The formalism provides: (i) Signal place partition ($\Psi \subseteq P$) distinguishing regulatory information channels from metabolic places, (ii) Signal flow arc semantics with consumption, (iii) Hierarchical layer assignment ($L: \Psi \to \mathbb{N}$) formalizing multi-scale organization, (iv) Preemption mechanism where higher-layer signal consumption collapses decision spaces at lower layers. This bridges the fundamental gap between metabolic networks (mass transfer) and regulatory networks (information flow) within unified Bio-PN formalism.

\noindent\textbf{(2) Proof-of-Concept Application}---We demonstrated the formalism's applicability through \textit{Bacillus subtilis} sporulation, a canonical ATP-gated developmental decision. Signal hierarchy analysis computes ATP commitment threshold at 2.38 mM, matching experimental measurements (2.21 mM, 7\% error). This quantitative prediction is inexpressible in classical Bio-PN lacking consumption semantics for regulatory signals. Similar analytical results were obtained for other biological systems (lambda phage lysis-lysogeny, MAPK signaling cascades), confirming the approach's generalizability. The application illustrates: \textit{Threshold determination}---Computing energy concentrations triggering pathway commitment, \textit{Basin geometry quantification}---Reachability analysis quantifying cell fate reliability and attractor convergence, \textit{Decision finality analysis}---Formal distinction between irreversible commitment (signal consumption) and reversible regulation (test arcs). This illustrative example establishes proof-of-concept for the formalism's predictive capability and demonstrates advantages over classical Bio-PN approaches.

\noindent Together, these contributions provide theoretical foundation for a new class of analyses about commitment reliability, basin geometry, and decision finality that lack formal semantics in classical Bio-PN theory.

\subsection{Organization}

Section~2 reviews classical Petri net theory, biological Petri nets, and identifies the expressiveness gap. Section~3 develops signal hierarchy theory with formal semantics. Section~4 presents the unified formalism integrating signal hierarchy with Bio-PN. Section~5 demonstrates applicability through the \textit{B. subtilis} sporulation example. Section~6 discusses theoretical significance and enabled analytical capabilities. Section~7 concludes with broader implications for quantitative systems biology.
